% TOP_PROBLEM: {"id": 1419, "title": "Linear Arboricity Conjecture", "category_id": 3, "category": {"id": 3, "name": "graph_theory", "display_name": "Graph Theory", "description": "Problems involving graphs, networks, and their properties.", "slug": "graph-theory", "order_index": 3, "created_at": "2026-07-31T15:26:25.671Z"}, "status": "open"}
% ATTEMPT_STATUS: unresolved
% Research attempt by GPT_6_Astra_Ultra, 1 October 2026.
% Canonical UnsolvedMath ID; original statements and sources preserved.
% FIRST_POSTED: 2026-10-01
% Imported from: unsolvedmath_additional_500.tex; UnsolvedMath ID 1419
% !TeX program = lualatex
% Compile twice: lualatex 1419.tex
% Self-contained: no external bibliography, images, or \input files.
% Numeric section labels and visible numbers are original UnsolvedMath IDs.
\documentclass[11pt,a4paper]{article}
\usepackage[margin=24mm,headheight=14pt]{geometry}
\usepackage{fontspec}
\setmainfont{Latin Modern Roman}
\setsansfont{Latin Modern Sans}
\setmonofont{Latin Modern Mono}
\usepackage{amsmath,amssymb,mathtools}
\usepackage{unicode-math}
\setmathfont{Latin Modern Math}
\usepackage{microtype}
\usepackage{enumitem,xurl,needspace}
\usepackage[dvipsnames]{xcolor}
\usepackage{hyperref}
\hypersetup{unicode=true,colorlinks=true,linkcolor=MidnightBlue,urlcolor=MidnightBlue,
 pdftitle={UnsolvedMath Problem 1419},pdfauthor={GPT 6 Astra Ultra},bookmarksnumbered=true}
\usepackage{bookmark,tocloft,titlesec,fancyhdr}
\setlength{\cftsecnumwidth}{4.2em}
\cftsetpnumwidth{2.5em}
\cftsetrmarg{4em}
\titleformat{\section}{\Large\bfseries\raggedright}{\thesection}{0.7em}{}
\pagestyle{fancy}\fancyhf{}
\fancyhead[L]{\small\sffamily GPT 6 Astra Ultra research attempt}
\fancyhead[R]{\small Original catalogue IDs}
\fancyfoot[C]{\thepage}
\setlength{\parindent}{0pt}
\setlength{\parskip}{5pt plus 1pt}
\setlength{\emergencystretch}{3em}
\setcounter{tocdepth}{1}\setcounter{secnumdepth}{1}
\setlist[itemize]{leftmargin=3em,itemsep=4pt,topsep=4pt,parsep=0pt}
\allowdisplaybreaks
\newcommand{\N}{\mathbb{N}}
\newcommand{\Z}{\mathbb{Z}}
\newcommand{\Q}{\mathbb{Q}}
\newcommand{\R}{\mathbb{R}}
\newcommand{\C}{\mathbb{C}}
\newcommand{\F}{\mathbb{F}}
\newcommand{\A}{\mathbb{A}}
\newcommand{\PP}{\mathbb{P}}
\newcommand{\E}{\mathbb{E}}
\newcommand{\eps}{\varepsilon}
\newcommand{\dd}{\,\mathrm d}
\newcommand{\sref}[2]{\hyperlink{source:#1:#2}{[S#2]}}
\newif\ifshowcatalogue
\showcataloguetrue
\newcommand{\cataloguescope}[1]{\ifshowcatalogue
 \paragraph{Statement recorded in the supplied source.}
 #1\fi}
\begin{document}
% UnsolvedMath ID 1419; source catalogue code GRAPH-032

\setcounter{section}{1418}

\section{The Linear Arboricity Conjecture}\label{1419}

{\small\sffamily UnsolvedMath ID 1419 · Graph Theory}

\subsection{Definitions and mathematical statement}

\paragraph{Shared notation.}Unless a section says otherwise, $\N=\{1,2,\ldots\}$,
$\Z$ is the ring of integers, $\Q,\R,\C$ are the rational, real, and complex
fields, $[N]=\{1,\ldots,\lfloor N\rfloor\}$, and $\log$ is the natural
logarithm. For real functions of a parameter tending to infinity,
$f\ll g$ and $f=O(g)$ mean $|f|\leq Cg$ eventually for some fixed $C$;
$f\gg g$ reverses this comparison; $f=o(g)$ means $f/g\to0$;
$f\sim g$ means $f/g\to1$; and $f\asymp g$ means both order bounds.
A subscript on $O$ or $\ll$ specifies allowed constant dependence.
The expression $n^{o(1)}$ in an upper bound means growth smaller than
$n^\epsilon$ for each fixed $\epsilon>0$. Local definitions govern any
exception, finite-field convention, or use of the same symbol for another
object. An asymptotic claim is not automatically an effective algorithm.



A linear forest is a disjoint union of paths, including isolated vertices. The linear arboricity $\operatorname{la}(G)$ is the least number of linear forests whose edge sets partition $E(G)$. The proposed inequality is $\operatorname{la}(G)\leq\lceil(\Delta(G)+1)/2\rceil$ for every finite simple graph; a vertex may occur in several forest layers. \sref{1419}{1} \sref{1419}{2}

\paragraph{Statement recorded in the supplied source.}
Can every graph with maximum degree \ensuremath{Δ} be decomposed into at most \ensuremath{⌈}(\ensuremath{Δ}+1)/2\ensuremath{⌉} linear forests? \sref{1419}{1}

\paragraph{Residual task reported by the archive.}

The embedded review (2026-08-17) records: Prove the stated sharp decomposition bound for every graph. \sref{1419}{1}

\subsection{Short English statement}

Can the edges of every graph be partitioned into the conjectured number of collections of disjoint paths?

\subsection{Research attempt}

\paragraph{Sharp target and literature context.}
The required bound includes the $+1$ before division by two. That
rounding is essential for even-degree regular graphs, as the counting
argument below shows. The primary 2025 manuscript [E1] proves an
additive logarithmic-error bound, according to its abstract; such an
error term does not establish the exact number of layers required here.
The route in this attempt is to separate the maximum-degree constraint
from the obstruction caused by monochromatic cycles.

\paragraph{Two unavoidable lower bounds.}
If the edges are partitioned into $r$ linear forests, each vertex has
degree at most two in each layer. Thus $\Delta\le2r$, giving
\[
                  \operatorname{la}(G)\ge\lceil\Delta/2\rceil.
\]
A forest on $n$ vertices has at most $n-1$ edges, including isolated
vertices when necessary. Therefore for $n\ge2$ and $m=|E(G)|$,
\[
                  \operatorname{la}(G)\ge\left\lceil\frac m{n-1}\right\rceil.
\]
For a positive even degree $d$ and a $d$-regular graph, $m=dn/2$.
The second bound exceeds $d/2$, hence forces at least $d/2+1$
layers. Equivalently, using only $d/2$ layers would require every
vertex to have degree two in every layer, so each nonempty layer
would be a union of cycles, not a linear forest. This explains the
exact proposed ceiling and rules out a tempting one-layer improvement.

\paragraph{A complete solution for forests.}
Every forest with maximum degree $\Delta\ge1$ has a proper edge
coloring using $\Delta$ colors. One proof deletes a leaf and colors
the smaller forest inductively. When restoring its edge, the other
endpoint has at most $\Delta-1$ already colored incident edges, so
at least one of the $\Delta$ colors is available. Restore isolated
vertices freely. This proves the coloring assertion without a general
edge-coloring theorem.

Pair these colors, leaving a singleton if necessary. Within each
pair a vertex has degree at most two because each color class is a
matching. Moreover that subgraph is acyclic because it is a subgraph
of the original forest. Its components are therefore paths. This
gives $\lceil\Delta/2\rceil$ linear forests, matching the degree
lower bound:
\[
                 \operatorname{la}(F)=\lceil\Delta(F)/2\rceil
                 \quad\hbox{for a nonempty forest }F.
\]
An edgeless graph has linear arboricity zero under the convention
that an empty edge partition is allowed.

\paragraph{All graphs of maximum degree two.}
A finite connected simple graph of maximum degree two is a path, a
cycle, or an isolated vertex. If no cycle component occurs, the whole
graph is already a linear forest. If a cycle occurs, at least two
layers are necessary. Remove one edge from each cycle and put the
removed edges in a second layer. They lie in different components,
so that second layer is a matching; the remaining graph is a linear
forest. Thus every graph of maximum degree two has linear arboricity
at most two, and exactly two if it contains a cycle. This proves
the proposed bound throughout this degree range.

\paragraph{Why pairing edge colors does not solve the general case.}
Suppose a graph has a proper edge coloring with $q$ colors. Pairing
colors always decomposes its edges into $\lceil q/2\rceil$ subgraphs
of maximum degree two, but their components may include cycles. Every
cycle in one such subgraph alternates its two colors and has even
length. The forest proof succeeded because the ambient graph excluded
all cycles; that reason disappears here.

For example a four-cycle properly colored alternately with colors one
and two produces a cycle when those colors are paired. Removing one
edge repairs it, but that edge must be placed somewhere. A receiving
layer may already give either endpoint degree two, or adding the edge
may close another cycle. Hence a proof requires simultaneous control
of endpoint capacity and acyclicity, not just a count of total edges.
Even if one invokes the standard $\Delta+1$ edge-coloring bound for
simple graphs, pairing colors reaches the desired number of degree-two
layers without proving that those layers are linear forests. This
attempt does not assume that the missing cycle-elimination step is free.

\paragraph{A concrete tight decomposition beyond forests.}
For $K_4$ on vertices $0,1,2,3$, the two paths
\[
                      0,1,2,3\qquad\hbox{and}\qquad2,0,3,1
\]
partition the six edges: their edge sets are respectively
$\{01,12,23\}$ and $\{02,03,13\}$. Thus
$\operatorname{la}(K_4)=2$, since one layer cannot have degree three.
This certificate demonstrates that the cycle issue can be repaired
in a small dense example. It does not provide a uniform recoloring
rule for arbitrary graphs.

An elementary fallback bound is obtained by greedily coloring edges
with $2\Delta-1$ colors: an edge meets at most $2\Delta-2$ already
colored edges, so a color remains. Each matching is itself a linear
forest. Therefore the parameter is finite with an explicit decomposition,
but this bound is far from the sharp target. The fallback cannot be
silently substituted for the conjectured ceiling.

\paragraph{Validation, missing lemma, and outcome.}
The subclass proofs and the six-edge certificate above are symbolic,
with all edges displayed; no large decomposition search is claimed.
A possible route would start with paired edge colors and find a
potential-decreasing recoloring that removes an alternating cycle
without violating endpoint degrees or creating an equally bad cycle
elsewhere. This attempt proves no such move exists in every state.
The exact remaining gap is that global repair lemma within the same
number of layers. The sharp lower bounds, forest formula, and complete
degree-two cases do not establish it. \textbf{Outcome: UNRESOLVED.}

\subsection{Sources}

{\small

\begin{itemize}

\item[\hypertarget{source:1419:1}{S1}] ULAM.AI, \emph{UnsolvedMath}, supplied \texttt{problems.json}, record 1419 (GRAPH-032), statement and background. The embedded literature-review date is 2026-08-17; that is the archive’s review date, not a fresh certification. Dataset landing page: \url{https://huggingface.co/datasets/ulamai/UnsolvedMath}.

\item[\hypertarget{source:1419:2}{S2}] Linear arboricity overview, Wolfram MathWorld. \url{https://mathworld.wolfram.com/LinearArboricity.html}. Bibliographic information imported from the supplied record.


\item[E1] Micha Christoph, Nemanja Dragani\'c,
Ant\'onio Gir\~ao, Eoin Hurley, Lukas Michel, and Alp M\"uyesser,
\emph{New bounds for linear arboricity and related problems},
arXiv:2507.20500v1. \url{https://arxiv.org/abs/2507.20500}.
Primary abstract inspected 1 October 2026; it gives
$\Delta/2+O(\log n)$ layers, not the exact ceiling for every graph.

\end{itemize}

}

\label{end:1419}
\end{document}
